%%%%%%%%%%%%%%%%%% USAGE INSTRUCTIONS %%%%%%%%%%%%%%%%%%
% - Compile using LuaLaTeX and biber, unless there is a particular reason not to. Do not use the older LaTex/PDFLaTeX or BibTeX (the fonts won't work correctly.)
% - Font and the report 'year' must be specified when all \documentclass or the template won't work correctly. (There's no error checking/default cases!)
% - Options for fonts are: calibri, times, palatino, garamond, arial, tahoma, verdana, trebuchet. Note however that not all these are installed on Overleaf, you need to install to your project 
% - Options for 'year' are: first, second, thesis.
% - If not a thesis, should probably remove the COVID-19 impact statement page
% - As many further packages as wanted can be loaded. Below are just an example set. Note that template itself loads a number of packages, including hyperref.
% - References are handed using biblatex.
% - Don't need to include a \uomdeclarations unless this is a thesis
% - Note there is more documentation in the header of uom_these_casson.cls file if you need more help
% - Link to the presentation of theses policy: http://www.regulations.manchester.ac.uk/pgr-presentation-theses/



%%%%%%%%%%%%%%%%%% META DATA SETUP %%%%%%%%%%%%%%%%%%
% This is where the document title and author are set. Other details for the title page are set later
\begin{filecontents*}{\jobname.xmpdata}
  \Title{A data reduction algorithm incorporating a low power continuous wavelet transform for use in wearable electroencephalography systems} 
  \Author{Alexander J. Casson}
  \Language{en-GB}
  \Copyrighted{True}
  % More meta-data fielda can be added here if wanted, see https://ctan.org/pkg/pdfx?lang=en for fields
\end{filecontents*}



%%%%%%%%%%%%%%%%%% DOCUMENT SETUP %%%%%%%%%%%%%%%%%%
\documentclass[times,thesis]{uom_thesis_casson} % See above for font options Year can be: first, second or thesis.



%%%%%%%%%%%%%%%%%% PACKAGES AND COMMANDS %%%%%%%%%%%%%%%%%%

% Packages - some useful examples
\usepackage{graphicx,psfrag,color} % for postscript graphics files
  \graphicspath{ {./images/} }
\usepackage{amsmath}               % assumes amsmath package installed
  \allowdisplaybreaks[1]           % allow eqnarrays to break across pages
\usepackage{amssymb}               % assumes amsmath package installed 
\usepackage{url}                   % format hyperlinks correctly
\usepackage{rotating}              % allow portrait figures and tables
\usepackage{multirow}              % allows merging of rows in tables
\usepackage{lscape}                % allows pages to be typeset in landscape mode
\usepackage{tabularx}              % allows fixed width tables
\usepackage{verbatim}              % enhanced version of built-in verbatim environment
\usepackage{footnote}              % allows more control over footnote environments
\usepackage{float}                 % allows H option on floats to force here placement
\usepackage{booktabs}              % improve table line spacing
\usepackage[base]{babel}           % required for lisum package
\usepackage{lipsum}                % for adding dummy text here
\usepackage{subcaption}            % for multiple sub-figures in a single float
\usepackage{siunitx}               % add SI units
% Add your packages here


% Custom commands
\newcommand{\degree}{\ensuremath{^\circ}}
\newcommand{\sus}[1]{$^{\mbox{\scriptsize #1}}$} % superscript in text (e.g. 1st can be 1\sus{st})
\newcommand{\sub}[1]{$_{\mbox{\scriptsize #1}}$} % subscript in text
\newcommand{\chap}[1]{Chapter~\ref{#1}}
\newcommand{\sect}[1]{Section~\ref{#1}}
\newcommand{\fig}[1]{Fig.~\ref{#1}}
\newcommand{\tab}[1]{Table~\ref{#1}}
\newcommand{\equ}[1]{(\ref{#1})}
\newcommand{\appx}[1]{Appendix~\ref{#1}}
% Add your commands here



%%%%%%%%%%%%%%%%%% REFERENCES SETUP %%%%%%%%%%%%%%%%%%

% Setup your references here. Change the reference style here if wanted
\usepackage[style=ieee,backend=biber,backref=true,hyperref=auto]{biblatex}
% Note backref=true adds a page number (and hyperlink) to each reference so you can easily go back from the references to the main document. You may prefer backref=false if you need to stick strictly to a given reference style


% Fixes which can't be applied in the .cls file
\DefineBibliographyStrings{english}{backrefpage = {cited on p\adddot},  backrefpages = {cited on pp\adddot}}
  \renewcommand*{\bibfont}{\large}


% Add more .bib files here if wanted
\addbibresource{references.bib}



%%%%%%%%%%%%%%%%%% START DOCUMENT %%%%%%%%%%%%%%%%%%
\begin{document}



%%%%%%%%%%%%%%%%%% COVID-19 impact statement %%%%%%%%%%%%%%%%%%
\begin{uomcovid} % policy asks for statement to be before the title page
  % Policy asks for this to be removed from the final thesis post-examination, which messes up the page numbers somewhat. Done here so that the final post-examination thesis is correct. In the pre-examination thesis the numbers displayed on the page will be one lower than the number displayed by the PDF reader. Remove this page if not wanted/needed, and/or is a first or second year report
  Optional. Delete this section if not needed. This is COVID-19 impact text.
  
  \lipsum[1-3] % generate dummy text for here
\end{uomcovid}% \clearpage is added automatically



%%%%%%%%%%%%%%%%%% TITLE PAGE %%%%%%%%%%%%%%%%%%

% Title and author are automatically taken from the document meta-data defined above
\makeatletter
\title{\xmp@Title}
\author{\xmp@Author}
\makeatother

% Set the below yourself
\faculty{Science and Engineering}                  % "Faculty of" is added automatically
\department{Department of Electrical and Electronic Engineering} % regulations allow School, Division, or Department to be put here
\submitdate{2022}                                  % regulations ask only for the year, not month
\wordcount{1000}		                           % use \wordcount{} to set the count, \thewordcount to print in the text
\maketitle



%%%%%%%%%%%%%%%%%% LISTS OF CONTENT %%%%%%%%%%%%%%%%%%

% Probably don't need all of these unless final thesis
\uomtoc % contents 
\uomlof % figures
\uomlot % tables
\begin{uomlop} % list of publications.
  % Can use biblatex to \printbibliography[heading=none] to populate automatically, or can add a custom list via the tocloft package, but probably easier to just type in unless you have lots!
  Publications go here.
\end{uomlop}
\begin{uomterms}
  Enter terms and abbreviations as table or similar
  % Add list of terms and abbreviations by hand if wanted. Is no formal requirement to have one. There are lots of packages to help with this if you don't want to do by hand
\end{uomterms}



%%%%%%%%%%%%%%%%%% ABSTRACT %%%%%%%%%%%%%%%%%%
\begin{abstract} % put abstract here. Limit is 1 page.
  This is abstract text. 
  
  \lipsum[1-2]
\end{abstract}%
\clearpage



%%%%%%%%%%%%%%%%%% LAY ABSTRACT %%%%%%%%%%%%%%%%%%
\begin{uomlay} % put lay abstract here. Limit is 1 page. Not compulsory
  This is lay abstract text. 
  
  \lipsum[1-2]
\end{uomlay}



%%%%%%%%%%%%%%%%%% DECLARATIONS %%%%%%%%%%%%%%%%%%
\uomdeclarations % Don't need unless final thesis. No options are needed. Having this command will add the required declarations



%%%%%%%%%%%%%%%%%% LIST OF THESIS REVISIONS %%%%%%%%%%%%%%%%%%
\begin{uomlotr} % Only required for resubmitted theses
Put list of revisions here. Only required for resubmitted theses. Delete if not needed
\end{uomlotr} 



%%%%%%%%%%%%%%%%%% ACKNOWLEDGEMENTS %%%%%%%%%%%%%%%%%%
\begin{uomacknowledgements} % probably don't need unless final thesis
Acknowledgements go here.
\end{uomacknowledgements}



%%%%%%%%%%%%%%%%%% ABOUT THE AUTHOR %%%%%%%%%%%%%%%%%%
\begin{uomauthor} % Optional
If desired, a brief statement for External Examiners giving the candidate’s degree(s) and research experience, even if the latter consists only of the work done for this thesis.
\end{uomauthor} 



%%%%%%%%%%%%%%%%%% CHAPTER 1 %%%%%%%%%%%%%%%%%%
\chapter{Introduction} % probably better to use \input{} for each chapter to draw content from other files rather than having everything in one big file
  
  \lipsum[1-5] % generate dummy text for here

\chapter{Literature review}

  \section{Introduction}
  \lipsum[1] 
  
  \section{Content}
  \lipsum[1-2] \cite{ref:jCAS09,ref:jCAS09a,ref:jCAS10} \lipsum[3-5] Below is an example table. See \tab{table:example_table} for more details.
  \begin{table}
    \centering
    \caption{Example table.}
    \label{table:example_table}
    \begin{tabular}{ccccc}
      \toprule
      \multirow{2}{*}{Participant} & \multicolumn{2}{c}{Number / \%} & \multicolumn{2}{c}{Duration / \%} \\
                                   & Prime dresses & Non-prime dresses & Prime dresses & Non-prime dresses \\
	  \toprule
      1  & 33.33 & 33.91 & 20.83 & 18.42 \\
      2  & 13.04 & 17.50 & 04.93 & 07.62 \\
      3  & 22.73 & 20.10 & 13.00 & 08.20 \\
      4  & 31.34 & 21.88 & 10.57 & 11.09 \\
      5  & 08.47 & 19.32 & 03.04 & 09.73 \\
      \hline
      Mean & 16.4 & 16.5 & 07.8 & 07.5 \\
      Standard deviation & 09.7 & 06.6 & 05.4 & 03.3 \\
      \bottomrule
    \end{tabular}
  \end{table}  
  
  This is an example equation. See \equ{equ:example} for info.
  \begin{equation}
      a^{2} + b^{2} = c^{2} \label{equ:example}
  \end{equation}
  
  \section{Summary}
  \lipsum[6]
  
  
  
%%%%%%%%%%%%%%%%%% CHAPTER 2 %%%%%%%%%%%%%%%%%%
\chapter{Really good work}

  \section{Introduction}
  \lipsum[1] 
  
  \section{Content}
    \subsection{Introduction}
	\lipsum[1]
	
	\subsection{Detail}
    \lipsum[1-2] \cite{ref:jCAS10,ref:jCAS09,ref:jCAS09a} \lipsum[3-5] This is an example figure. See \fig{fig:uom_logo} for more details. 
	\begin{figure}
      \centering
      \includegraphics[width=0.3\textwidth]{uom_logo.pdf}
      \caption{Example figure.}
      \label{fig:uom_logo}
    \end{figure} 
    
	\subsection{More detail}
	\lipsum[1-2] \cite{ref:jCAS10,ref:jCAS09,ref:jCAS09a} \lipsum[3-5]
	
	\subsection{Summary}
	\lipsum[1]
  
  \section{Summary}
  \lipsum[6]
 



%%%%%%%%%%%%%%%%%% CONCLUSIONS %%%%%%%%%%%%%%%%%%
\chapter{Conclusions}



%%%%%%%%%%%%%%%%%% REFERENCES %%%%%%%%%%%%%%%%%%
\printbibliography[title={References},heading=bibintoc] % a single list of references for the whole thesis



%%%%%%%%%%%%%%%%%% APPENDICES %%%%%%%%%%%%%%%%%%
\begin{uomappendix} 
  \chapter{First appendix}
    \section{Section in Appendix}
    \lipsum[1-6]
\end{uomappendix}


%%%%%%%%%%%%%%%%%% END MATTER %%%%%%%%%%%%%%%%%%
\end{document}